\documentclass[10pt,a4paper]{amsart}
\usepackage[margin=19mm]{geometry}
\usepackage{amsmath,amssymb,mathtools}
\usepackage[scr=rsfs,cal=euler]{mathalfa}
\usepackage{xfrac}
\usepackage[unicode,hidelinks,bookmarks=false]{hyperref}
\newcommand{\ie}{i.e.\ }
\newcommand{\cf}{cf.\ }
\newcommand{\resp}{respectively\ }
\newcommand{\Subsection}{\subsection}
\providecommand{\nobreakdash}{\nobreak\hbox{-}}
\setlength{\emergencystretch}{2em}
\multlinegap0pt
\usepackage{amssymb}
\usepackage{tikz}
\usepackage[shortlabels]{enumitem}
\newtheorem{defn}{Definition}
\newtheorem{thm}{Theorem}
\newtheorem{prop}{Proposition}
\newtheorem{lem}{Lemma}
\title{Maximal Center of Distances of Finite Ultrametric Spaces and Perfect Binary Trees}
\author{Oleksiy Dovgoshey, Navjay Singh Jaiswal, Surinder Pal Singh Kainth, Gursimar Kaur, Olga Rovenska}
\date{Source: arXiv:2609.11384v1 (CC BY 4.0)}
\begin{document}
\maketitle

\noindent\emph{Source and licence.} Maximal Center of Distances of Finite Ultrametric Spaces and Perfect Binary Trees. Version arXiv:2609.11384v1 (CC BY 4.0), distributed by arXiv under the Creative Commons Attribution 4.0 International licence, http://creativecommons.org/licenses/by/4.0/, as recorded in the arXiv licence metadata for this exact version and shown on the abstract page https://arxiv.org/abs/2609.11384v1. Full licence notice: \texttt{licenses/arxiv-2609.11384-CC-BY-4.0.txt}.\par\smallskip

\noindent\textit{Editorial scope.} Counted unit: the article's thm m0 (source0.tex lines 1886--1915). The article's complete original proof of it is delivered with it (source0.tex lines 1917--2070, one \texttt{proof} environment of 154 lines). No mathematics is written, repaired or re-derived: the statement and the proof are the article's own lines, in the article's own notation, and no retained line is substituted or re-flowed.  Retained uncounted as the source-local context the proof needs: lines 235--241; lines 799--805; lines 810--817; lines 903--908; lines 918--932; lines 1118--1121; lines 1224--1254; lines 1493--1515.  Nothing was omitted from any retained span, so the package carries no omission record.  No retained line contains a citation, so the wrapper carries no bibliography and no bibliography entry.  No retained line contains a figure environment, an image-inclusion command or drawing code, so no image is delivered and the package declares no asset dependency.  Editorial formatting only, in addition: an AMS \texttt{amsart} preamble that restates the article's own declarations and macros used by the retained lines, namely the environments \texttt{defn}, \texttt{thm}, \texttt{prop}, \texttt{lem} on separate counters, together with the standard packages those lines need.  The retained environments print in this document as Definition 1, Theorem 1, Proposition 1, Lemma 1 in order of appearance, whereas the article prints the same items with the numbering of its own full document.  The complete native source of the article as distributed in the arXiv e-print for this exact version is bundled unchanged as \texttt{sources/210015/source0.tex}.
\begin{defn}\label{zcgh572}
Let $(X,d)$ be a metric space and let $D(X)$ be the distance set of $(X,d)$,
\begin{equation*}
D(X) := \{ d(x,y) : x,y \in X \}.
\end{equation*}
The {\it center of distances} $C(X)$  of $(X,d)$ is the set of all $t\in D(X)$ such that for each $p\in X$  there exists $x\in X$ with $d(p,x)=t$.
\end{defn}

\begin{defn}\label{edrtaa}
Let $k\geq 2$ be an integer number.
A non-empty graph $G$ is called complete $k$-partite if its vertices can be divided into disjoint non-empty sets $X_1,\ldots,X_k$ so that there are no edges joining the vertices which belong to the same $X_i$, and, moreover, any two vertices from different $X_i$ and $X_j$ are adjacent. In this case we write
\(
G=G[X_1,\ldots,X_k].
\)
\end{defn}

\begin{defn}\label{okiu}
Let $(X,d)$ be a finite metric space with $|X|\geq 2$. The diametrical graph $G_X$ of $(X,d)$ is a graph such that $V(G)=X$ and, for all $u,v\in X$,
\[
\{u,v\}\in E(G_{X})
\iff
d(u,v)=\operatorname{diam}X.
\]
\end{defn}

\begin{thm}\label{rvqqbb}
Let $(X,d)$ and $(Y,\rho)$ be  finite ultrametric spaces. Then the representing trees
$T_X=T_X(r_X,
l_X)$ and $T_Y=T_Y(r_Y,l_Y)$ are isomorphic as labeled rooted trees if and only if
$(X,d)$ and $(Y,\rho)$ are isometric.
\end{thm}

\begin{prop}\label{effrt}
Let $(X,d)$ and $(Y,\rho)$ be finite ultrametric spaces with  representing trees $T_X=T_X(r_X,l_X)$ and $T_Y=T_Y(r_Y,l_Y)$, respectively. Then 
\(
T_X(r_X,l_X)
\) and \(
T_Y(r_Y,l_Y)
\)
are isomorphic as labeled rooted trees if and only if
\(
T_X(l_X)
\) and
\(T_Y(l_Y)
\)
are isomorphic as labeled trees.
\end{prop}

\begin{equation}
    \label{mkkggd}
D_p(X) := \{\, d(p,x) : x \in X \,\}.
\end{equation}

\begin{lem} \label{UT2Lemma3}
Let $(X,d)$ be a ${\bf MCD}$-space and let $n \geq 1$ be an integer number such that 
\begin{equation}
    \label{fl}
|X|=2^n
\quad\text{and}\quad
|C(X)|=n+1.
\end{equation}
\begin{enumerate}
\item[(i)]  The representing tree
\(
T_{ X}=T_{ X}(r_{ X},l_{ X})
\)
of $( X,d)$ is a binary tree.

\item[(ii)]  The equality
\begin{equation*}
\operatorname{lev}_{T_X(r_X)}(v)=n
\end{equation*}
holds for each leaf $v$ of $T_{ X}(r_X)$.

\item[(iii)] The equalities
\begin{equation}
    \label{o1}
D(X)=C(X)=\{\,l_X(v):v\in V(T_X)\,\}=D_p(X)
\end{equation}
hold for each \(p\in X\),
where $D( X)$ is the distance set of $( X,d)$ and
$D_p( X)$ is defined by formula~\eqref{mkkggd}.
\end{enumerate}
\end{lem}

\begin{lem}
    \label{lemlem}
Let $(X,d)$ be a ${\bf MCD}$-space, let $n\geq 1$ be an
integer number satisfying the equalities
\begin{equation}
    \label{ddaad}
|X|=2^n \quad\text{and}\quad |C(X)|=n+1,
\end{equation}
and let $G(X)$ be the diametrical graph of $(X,d)$. Then
$G(X)$ is complete bipartite with parts $X_1,$ $X_2$ such that
\begin{equation}
    \label{ddad2}
|X_1|=|X_2|=2^{n-1}
\quad\text{and}\quad
|C(X_1)|=|C(X_2)|=n,
\end{equation}
where \(C(X_1)\) and \(C(X_2)\) are the centers of distances of ultrametric spaces
\(
(X_1,d|_{X_1\times X_1})\) and \(
(X_2,d|_{X_2\times X_2}),
\)
respectively.
\end{lem}

\begin{thm}\label{m0}
Let $(X,d)$ and $(Y,\rho)$ be $\bf MCD$-spaces.
Then the following statements are equivalent.
\begin{enumerate}
\item $(X,d)$ and $(Y,\rho)$ are isometric.

\item The representing trees
\(
T_X=T_X(r_X,l_X)
\) and
\(T_Y=T_Y(r_Y,l_Y)
\)
are isomorphic as labeled rooted trees.

\item The trees $T_X(l_X)$ and $T_Y(l_Y)$ are isomorphic as labeled trees.

\item The equality
\begin{equation}
    \label{m2}
C(X)=C(Y)
\end{equation}
holds.

\item The equality
\begin{equation*}
D(X)=D(Y)
\end{equation*}
holds.
\end{enumerate}
\end{thm}

\begin{proof} The equivalence 
$(i)\iff (ii)$ is valid by Theorem~\ref{rvqqbb}.

The validity of the equivalence $(ii) \iff (iii)$ follows from Proposition~\ref{effrt}.

Statement~$(iii)$ of Lemma~\ref{UT2Lemma3} implies the validity of 
\(
(iv)\iff(v).
\)
 It follows from Definition \ref{zcgh572} that the centers of distances
of isometric metric spaces are the same. Hence the implication
$(i) \Rightarrow (iv)$ is valid also. 

Thus to complete the proof  it suffices
to prove the validity of the implication
$(iv) \Rightarrow (i).$


Let statement $(iv)$ hold.
Then equality \eqref{m2} implies the existence of an integer $n \ge 0$ such that
\begin{equation}
    \label{v1-s}
|C(X)| = |C(Y)| = n+1. 
\end{equation}



We will prove that $(X,d)$ and $(Y,\rho)$ are isometric by induction.
If $n=0$ then we have 
\begin{equation*}
|X| = |Y| = 1,
\end{equation*}
since the definition of ${\bf MCD}$-spaces and \eqref{v1-s} imply 
\begin{equation*}
|X| = |Y| = 2^n=1.
\end{equation*}
Since all
single-point metric spaces are isometric, statement $(i)$ is valid for  $n=0.$


Let us consider the case $n\geq 1$. Let $n_0 \geq 1$ be an integer number and let statement $(i)$ hold for all ${\bf MCD}$-spaces $(X,d)$, $(Y,\rho)$ satisfying \eqref{v1-s} with $n<n_0$. 

By Lemma \ref{lemlem} the diametrical graphs
$G_X$ and $G_Y$ are complete bipartite. Let $X_1,$ $X_2$ be the parts of
$G_X$, and $Y_1,$ $Y_2$ be the parts of $G_Y$. By Lemma \ref{lemlem} the
ultrametric spaces
\(
(X_1,d|_{X_1\times X_1}),\)
\((X_2,d|_{X_2\times X_2}),\)
\((Y_1,\rho|_{Y_1\times Y_1}),\)
\((Y_2,\rho|_{Y_2\times Y_2})
\)
belong to ${\bf MCD}$ and the equalities
\begin{equation*}
    |X_i|=|Y_i|=2^{n_0-1},
\end{equation*}
\begin{equation}
    \label{ddaad4}
    |C(X_i)|=|C(Y_i)|=n_0
\end{equation}
hold for $i=1,2$ if \eqref{v1-s} holds with $n=n_0$.
The induction hypothesis and equalities \eqref{ddaad4} imply that all spaces
\(
(X_1,d|_{X_1\times X_1}),
\)
\(
(X_2,d|_{X_2\times X_2}),
\)
\(
(Y_1,\rho|_{Y_1\times Y_1})
\) and \(
(Y_2,\rho|_{Y_2\times Y_2})
\)
are isometric. Let
\(
\Phi_1:X_1\to Y_1
\)
and
\(
\Phi_2:X_2\to Y_2
\)
be the corresponding isometries and let a mapping 
\(
\Phi:X\to Y
\)
be defined as
\begin{equation}
    \label{ddaad5}
\Phi(p):=
\begin{cases}
\Phi_1(p), & \text{if }p\in X_1,\\
\Phi_2(p), & \text{if }p\in X_2.
\end{cases}
\end{equation}
The mapping $\Phi:X\to Y$ is an isometry of $(X,d)$ and $(Y,\rho)$ if and only if the equality
\begin{equation}
    \label{ddaad6}
d(p,q)=\rho(\Phi(p),\Phi(q))
\end{equation}
holds for all $p,q\in X$.
Let us prove equality \eqref{ddaad6}. If
\(
p,q\in X_1 \) or
\( p,q\in X_2,
\)
equality \eqref{ddaad6} follows from \eqref{ddaad5}, since
\(
\Phi_1:X_1\to Y_1
\) and \(
\Phi_2:X_2\to Y_2
\)
are isometries. Suppose that $p$, $q$ belong to distinct parts $X_1,$ $X_2$
of $G_X$. Then, using  equality \eqref{ddaad5}, we see that
\(
\Phi(p)\), \(\Phi(q)
\)
belong to distinct parts $Y_1,$ $Y_2$ of $G_Y$. Now,
Definitions \ref{edrtaa}, \ref{okiu} imply the equalities 
\begin{equation}
    \label{ttr1}
d(p,q)=\operatorname{diam} X,
\end{equation}
\begin{equation}
    \label{ttr2}
\rho(\Phi(p),\Phi(q))
=\operatorname{diam} Y. 
\end{equation}
By Lemma \ref{UT2Lemma3} the equalities
\begin{equation*}
C(X)=D(X),\quad C(Y)=D(Y) 
\end{equation*}
hold and, consequently, we have
\begin{equation*}
    D(X)=D(Y)
\end{equation*}
by equality \eqref{m2}.
The last equality and the definition of diameter of metric spaces imply
\[
\operatorname{diam} X
=
\max\{t:t\in D(X)\}
=
\max\{t:t\in D(Y)\}
=
\operatorname{diam} Y.
\]
Equality \eqref{ddaad6} follows from \eqref{ttr1}, \eqref{ttr2}.

Thus the implication
\(
(iv)\Rightarrow(i)
\)
is valid.
\end{proof}

\end{document}
